\section{Benchmark Contract}
\label{sec:training}

The experiments follow the same decomposition as the theorem: transport, symbol, residual, packetization, and refinement are measured separately whenever the data support it.  The benchmark uses a unified scenario registry over synthetic surrogates, cached TCNO PDE families, imported mature reference rows, and optional sequence-valued rollout paths.  One-step approximation rows are kept separate from nonlinear rollout diagnostics.

\paragraph{Model split.}
All theorem-facing claims attach to \MiNO-Core: analytic packet tokenizer, bounded canonical stencil, local symbol branch, transported synthesis, transported high-frequency carrier, and explicit residual.  \MiNO-Plus keeps the same packet backbone and adds confidence-gated transport and a low-pass post-synthesis local refinement head.  Reported gains from \MiNO-Plus are empirical enlargements of the theorem-facing Core, while the carrier-bound mechanism test still uses the same transported packet primitive.

\paragraph{Staged runs.}
Stage~0 is a two-epoch smoke check. Stage~1 is a one-seed breadth pass over the synthetic and cached PDE slate. Stage~2 repeats the most important wave, diffusion, and Helmholtz rows over three seeds. Stage~3 is the higher-budget oscillatory stress path.  Runs are written as raw JSON plus aggregate CSV rows, and the empirical-closure runner uses stable row identifiers so interrupted campaigns resume with \texttt{--skip-existing}.

\paragraph{Objectives.}
The breadth and shortlist matrices use field loss,
\[
\mathcal L_{\mathrm{field}}=\|u_{\mathrm{pred}}-u_{\mathrm{target}}\|_2^2.
\]
The theorem-aligned mechanism tests add packet proxies,
\[
\mathcal L_{\mathrm{field}}
+\lambda_{\mathrm{tr}}\mathcal L_{\mathrm{tr}}
+\lambda_{\mathrm{sym}}\mathcal L_{\mathrm{sym}}
+\lambda_{\mathrm{can}}\mathcal L_{\mathrm{can}}
+\lambda_{\mathrm{ord}}\mathcal L_{\mathrm{ord}}
+\lambda_{\mathrm{pkt}}\mathcal L_{\mathrm{pkt}}
+\lambda_{\mathrm{hf}}\mathcal L_{\mathrm{hf}}
+\lambda_{\mathrm{flow}}\mathcal L_{\mathrm{flow}},
\]
where the terms measure transported packet geometry, packet-amplitude mismatch, canonical consistency, symbol-order behavior, direct packet-space reconstruction, high-frequency core reconstruction, and controlled-flow metadata displacement when a known packet drift is available.  The carrier-bound profile trains the core first, freezes local refinement during warmup, releases the refinement head only as a low-pass residual, and routes high-frequency reconstruction through transported packet metadata.

The executed carrier-bound weights are fixed in the runner rather than tuned per scenario:
\[
\begin{aligned}
\lambda_{\mathrm{tr}}&=0.20,&
\lambda_{\mathrm{sym}}&=0.05,&
\lambda_{\mathrm{can}}&=0.02,\\
\lambda_{\mathrm{ord}}&=0.01,&
\lambda_{\mathrm{pkt}}&=0.05,&
\lambda_{\mathrm{hf}}&=1.00,\\
\lambda_{\mathrm{flow}}&=2.00.&
\end{aligned}
\]
The mechanism profile uses \texttt{hamiltonian\_verlet} transport, \texttt{spectral\_order} symbol modulation, Gaussian/Gabor packets by default, sparse top-$K$ aggregation after neighbor selection, transported synthesis scale one, transported input-carrier scale one, token-level refinement scale zero, and skip/refinement low-pass cutoff $0.16$.

In the theorem-facing \texttt{spectral\_order} mode, the executable symbol action is multiplier-only: additive local shifts are disabled and belong instead to the residual or finite-landing terms.  The symbol-order proxy used in $\mathcal L_{\mathrm{ord}}$ is computed from the symbol output without detaching the packet values, so it is a training penalty as well as a logged diagnostic.  The generic \texttt{mlp} symbol branch may still use an affine scale--shift action in legacy empirical rows, but those shifts are not used to justify the finite \(S^m\)-proxy statement.

\paragraph{Learning certificate.}
The approximation theorem is an existence statement over the \MiNO\ hypothesis class; it does not claim that stochastic gradient descent reaches the global optimum.  The connection to supervised training is the standard excess-risk certificate below.  Let $\mathcal D$ be a distribution over input--target pairs $(u,Tu)$ and let
\[
\mathcal R_{\mathrm{field}}(\theta)
:=
\mathbb E_{(u,Tu)\sim\mathcal D}
\|M_\theta(u)-Tu\|_{L^2}^2 .
\]
Let the theorem-aligned training objective be
\[
\mathcal J(\theta)
:=
\mathcal R_{\mathrm{field}}(\theta)
+
\mathcal P_{\mathrm{proxy}}(\theta),
\qquad
\mathcal P_{\mathrm{proxy}}(\theta)\ge0,
\]
where $\mathcal P_{\mathrm{proxy}}$ collects the nonnegative transport, symbol, canonical, packet, high-frequency, flow, and optional Egorov penalties used by the selected profile.

\begin{proposition}[Realizability-to-learning certificate]
\label{prop:learning-certificate}
Assume the hypotheses of Theorem~\ref{thm:wave-mino-core} on the support of $\mathcal D$.  Suppose there exists a theorem-grade parameter $\theta^\star$ such that
\[
\|M_{\theta^\star}(u)-Tu\|_{L^2}
\le
e_{h,K,P}(u)
\quad
\text{for } \mathcal D\text{-almost every }u,
\]
and
\[
\mathcal P_{\mathrm{proxy}}(\theta^\star)\le \Pi_{\mathrm{proxy}} .
\]
Let $\widehat{\mathcal J}_n$ be the empirical version of $\mathcal J$ on $n$ training pairs.  If the trained parameter $\hat\theta$ satisfies the empirical optimization gap
\[
\widehat{\mathcal J}_n(\hat\theta)
\le
\inf_{\theta\in\mathcal H_{\MiNO}}\widehat{\mathcal J}_n(\theta)
+
\eta_{\mathrm{opt}},
\]
and the profile class has uniform generalization gap
\[
\sup_{\theta\in\mathcal H_{\MiNO}}
\bigl|\mathcal J(\theta)-\widehat{\mathcal J}_n(\theta)\bigr|
\le
\eta_{\mathrm{gen}},
\]
then
\[
\mathcal R_{\mathrm{field}}(\hat\theta)
\le
\mathbb E_{\mathcal D}\, e_{h,K,P}(u)^2
+
\Pi_{\mathrm{proxy}}
+
\eta_{\mathrm{opt}}
+
2\eta_{\mathrm{gen}} .
\]
\end{proposition}

\begin{proof}
Since $\mathcal P_{\mathrm{proxy}}\ge0$, $\mathcal R_{\mathrm{field}}(\hat\theta)\le \mathcal J(\hat\theta)$.  Uniform generalization gives
\[
\mathcal J(\hat\theta)
\le
\widehat{\mathcal J}_n(\hat\theta)+\eta_{\mathrm{gen}}
\le
\widehat{\mathcal J}_n(\theta^\star)+\eta_{\mathrm{opt}}+\eta_{\mathrm{gen}}
\le
\mathcal J(\theta^\star)+\eta_{\mathrm{opt}}+2\eta_{\mathrm{gen}}.
\]
The realizability hypothesis and proxy bound give
\[
\mathcal J(\theta^\star)
\le
\mathbb E_{\mathcal D} e_{h,K,P}(u)^2+\Pi_{\mathrm{proxy}},
\]
which proves the claim.
\end{proof}

Proposition~\ref{prop:learning-certificate} is intentionally modest.  It says that the microlocal approximation theorem becomes a supervised learning guarantee once the optimizer, sample size, and proxy weights deliver small $\eta_{\mathrm{opt}}$, $\eta_{\mathrm{gen}}$, and $\Pi_{\mathrm{proxy}}$.  The experiments therefore test not only expressivity but also whether the registered training protocol makes the theorem-grade mechanism reachable in finite budget.

\paragraph{Mechanism-identifiability protocol.}
\label{sec:transport-identifiability-campaign}
The completed \texttt{branch\_id\_v3} campaign is the main branch attribution test.  It uses \texttt{hamiltonian\_verlet} metadata transport, \texttt{gabor\_gaussian} tokenization, sparse top-$K$ aggregation after dense neighbor search, proxy losses, core-first warmup, low-pass residual refinement, transported synthesis, transported high-frequency input carrier, and same-refinement FNO/UNO/WNO controls.  Its ablations include full, no transport, identity transport, randomized metadata, no symbol, oracle transport where a known flow is available, and same-refinement baselines.

The ablation semantics are fixed in the runner.  \texttt{no\_transport} replaces the propagation block by a zero message and leaves metadata at the analysis centers, so the carrier path is still present but cannot move high-frequency content.  \texttt{identity\_transport} keeps the propagation/message block active while setting the learned displacement scale to zero.  \texttt{randomized\_metadata} permutes the packet metadata across tokens before propagation, preserving the metadata multiset while breaking coefficient--metadata association.  Same-refinement baselines attach the identical low-pass image-space refinement head, cutoff, route initialization, and training loss budget to the baseline prediction; they do not receive MiNO's packet tokenizer, transported synthesis, transported high-frequency carrier, or token-level packet refinement.

The interpretation rule is fixed: canonical transport is empirically credited only when transport corruption or metadata randomization degrades field or packet diagnostics, and when the same local refinement head does not by itself explain the microlocal branch effect.  The carrier-bound rows pass this rule on controlled wave probes: removing transport and randomizing metadata consistently worsens field error, while the same-refinement rows remain the absolute-accuracy controls.

The controlled-flow rows also use the known packet drift through $\mathcal L_{\mathrm{flow}}$.  This theorem-aligned supervision term tests whether a carrier-bound architecture can use a canonical metadata channel when the target flow is available.  The executable control campaign adds two fairness gates: a flow-off row, which removes $\mathcal L_{\mathrm{flow}}$ at fixed architecture, and a carrier-off transport row, which removes both transport and transported carrier paths.  These rows separate supervision dependence, carrier placement, and transport learning.

\paragraph{Egorov pilot diagnostics.}
Egorov consistency is evaluated as a diagnostic rather than as a default training loss.  Choose a fixed probe family $\mathcal A_{\mathrm{probe}}=\{A_{j,h}^{\mathrm{pkt}}\}_{j=1}^{J}$ of order-zero localization multipliers on the packet window.  For theorem-grade linear Core rows we evaluate $\rho_{\mathrm{Eg},h}^{\mathrm{pkt}}(j)$ from Section~\ref{sec:continuous}; for nonlinear \MiNO-Plus rows we evaluate the Jacobian version $\rho_{\mathrm{Eg},h}^{\mathrm{Jac}}(u,j)$.  The executable runner also exposes two cheaper controlled-flow probes: a self-intertwining probe comparing $A_{\mathrm{out}}M(u)$ with $M(A_{\mathrm{in}}u)$, and a target-calibrated probe comparing $M(A_{\mathrm{in}}u)$ with $A_{\mathrm{out}}v$ for the supervised target $v$.  The target-calibrated probe is the identifying one: it can detect a model that is self-consistent but transports to the wrong location.  If used as a training term, the full residual enters a theorem-aligned optional profile,
\[
\mathcal L_{\mathrm{field}}
+\lambda_{\mathrm{Eg}}
\frac1J\sum_{j=1}^{J}
\rho_{\mathrm{Eg},h}^{\mathrm{Jac}}(u,j)^2,
\]
with \texttt{egorov\_loss\_weight} reported separately.  The relevant ablations are \texttt{egorov\_on}, \texttt{egorov\_off}, \texttt{no\_transport}, and \texttt{no\_symbol}.  The diagnostic target is separation in the calculus metric: transport corruption should increase the pullback part of the residual, while symbol removal should increase the left/right symbol-coupling part.

As a minimal sanity check, we ran the finite-difference proxy on the bicharacteristic control row with three seeds, six epochs, and sample caps of 8/4/4.  The self-Jacobian proxy is not transport-identifying: it is $0.075$ for \texttt{full}, $0.067$ for \texttt{no\_symbol}, and $0.036$ for \texttt{no\_transport}, showing that a model can remain self-consistent under an incorrect transport.  The target-calibrated finite-difference Jacobian proxy is more informative: it is $0.402$ for \texttt{full}, $0.437$ for \texttt{no\_symbol}, and $0.848$ for \texttt{no\_transport}, matching the field-error ordering $0.426$, $0.465$, and $0.879$.  The target-calibrated proxy is therefore a useful calculus diagnostic; the main empirical verdict remains based on packet propagation and carrier-bound branch identifiability.

\paragraph{Cross-resolution protocol.}
\label{sec:cross-resolution-campaign}
The repository also exposes a \texttt{cross\_resolution\_wave} grid,
\[
64\to64,\qquad 64\to128,\qquad 96\to192,\qquad 128\to256,
\]
with fixed/scaled packet budgets and fixed/scaled stencil budgets.  It is retained as the executable counterpart of Theorem~\ref{thm:cross-resolution}, but it is not used as completed evidence in the current empirical verdict.
